%% ============================================================
\section{Method}\label{sec:method}
%% ============================================================

CRAFT applies one clause-processing pipeline to inference, SFT target
construction, and RL credit assignment.  Inference verifies the composition
of the filtered clause union, SFT serializes an accepted composition as one
response, and RL maps surviving contributions back to their proposing
rollouts.  Assertions and postconditions remain hidden throughout generation
and training and are restored only for final verification.
Appendix~\ref{app:formal-method} gives the complete formal definition.

\subsection{Target-Hidden Composable Inference}
\label{sec:method:core}
\label{sec:method:inference}

Pooling serves two purposes: independently valid clauses can jointly prove a
target that none proves alone, and clauses from different responses can
support one another during inductive verification.  CRAFT therefore pools
the decomposed clauses before filtering, rather than requiring each response
to succeed independently.

For a masked loop \(\mathcal L\), let \(\operatorname{parts}(y)\) be the
ordered invariant clauses extracted from response \(y\).  Let
\(\operatorname{Admit}(y)=\operatorname{Unique}(
\operatorname{First}_{m_{\max}}(\operatorname{parts}(y)))\), where
\(\operatorname{First}_{m}\) keeps the first \(m\) clauses.  The value of
\(m_{\max}\) is reported in Table~\ref{tab:constants}.
Algorithm~\ref{alg:method-pipeline} defines the complete inference procedure.

\begin{algorithm}[H]
\caption{Rollout--decompose--pool--filter--compose}
\label{alg:method-pipeline}
\begin{algorithmic}
\alginput{task \(\mathcal T=(\mathcal L,\mathcal G)\), policy \(\pi\), number of responses \(n\)}
\algoutput{composed invariant \(I\), final verdict \(v\)}
\algline \(x\leftarrow\mathcal L\) by masking \(\mathcal G\) in \(\mathcal T\)\cmnt{hide target}\\
\algline \(y_{1:n}\sim\pi(\cdot\mid x)\)\cmnt{rollout}\\
\algline \(A_i\leftarrow\operatorname{Admit}(y_i)\) for \(i=1,\ldots,n\)
  \cmnt{decompose and cap}\\
\algline \(P\leftarrow\operatorname{Unique}(A_1\Vert\cdots\Vert A_n)\)
  \cmnt{pool in response order}\\
\algline \(C\leftarrow\operatorname{Filter}_{\mathcal L}(P)\)\cmnt{filter}\\
\algline \(I\leftarrow\bigwedge_{c\in C}c\)\cmnt{compose}\\
\algline \(v\leftarrow\operatorname{TaskVerify}(\mathcal T,C)\)\cmnt{restore target}\\
\algline \kw{return} \((I,v)\)\\
\end{algorithmic}
\end{algorithm}

\tightparagraph{Joint verification and deletion.}
\(\operatorname{Unique}\) removes duplicates conservatively, keeping the first
occurrence in response and clause order.  The filter receives this ordered
family \(D_0\).  For each \(c\in D\), the loop verifier checks
\(\mathsf{Init}\Rightarrow c\) and \(\{I(D)\wedge B\}S\{c\}\), returning
\(\mathsf{valid}\), \(\mathsf{timeout}\), or \(\mathsf{unknown}\).
With per-obligation budget \(\delta\), one deletion round is
\begin{equation}
 D_{j+1}=\left\langle c\in D_j\;\middle|\;
 \operatorname{LoopVerify}_{\mathcal L,\delta}(D_j)[c]=\mathsf{valid}
 \right\rangle.
 \label{eq:main-houdini-round}
\end{equation}
After each deletion, survivors are rechecked because they may have depended
on removed clauses.  Only a proved fixed point is returned; timeout and
unknown cause deletion, and exhausting the round budget returns an empty
sequence.  Under a sound verifier, the resulting conjunction is inductive.
Appendix~\ref{app:formal-houdini} gives the bounded operator, and
Appendix~\ref{app:composition-proofs} establishes its soundness and the
candidate-relative optimality of ideal Houdini.

\tightparagraph{Why pooling precedes filtering.}
Let \(\mathsf H_{\mathcal L}\) denote ideal Houdini with exact logical
decisions.  Filtering the union retains every clause obtainable by separately
filtering the responses:
\begin{equation}
 \bigcup_i\mathsf H_{\mathcal L}(A_i)
 \subseteq\mathsf H_{\mathcal L}\!\left(\bigcup_i A_i\right).
 \label{eq:main-pooling-inclusion}
\end{equation}
The inclusion can be strict when supporting clauses occur in different
responses; its derivation is in Appendix~\ref{app:composition-proofs}.
Finite prover budgets can break this inclusion, so measured compose@\(k\)
need not be monotonic (Appendix~\ref{app:houdini-order}).

Only after \(I\) is fixed do we restore the hidden targets \(\mathcal G\)
and run the final verification conditions at their original locations.  The
inference objective is therefore
\[
  \max_\pi\ \Pr_{y_{1:n}\sim\pi(\cdot\mid\mathcal L)}
  [\operatorname{TaskVerify}((\mathcal L,\mathcal G),C(y_{1:n}))].
\]
Its empirical estimate is compose@\(n\): success is judged on the final
composition, rather than on any complete response.  Both complementary
target proofs and cross-response inductive support can therefore contribute
to gains over pass@\(n\).

\subsection{Target-Independent Strength Signal}
\label{sec:method:sampler}

Inductive validity alone does not distinguish useful invariants from weak
ones such as \(\mathsf{true}\).  With targets hidden, training therefore
needs a target-independent signal of how much behavior a valid invariant
excludes.  We execute the masked loop to collect reachable loop-head states
\(\mathcal E\) and construct \emph{candidate negative traces}
\(\mathcal N\).  Local relational perturbations probe constraints among
variables, while continuations beyond witnessed exits probe boundary
constraints.  Appendix~\ref{app:negative-sampling} gives the complete sampler,
including cleaning, per-family budgets, and random fill of unused capacity.

For a verified clause set \(C\), the state evaluator returns
\(\operatorname{Eval}(c,s)\in\{\mathsf{true},\mathsf{false},\mathsf{unknown}\}\).
A negative trace is rejected only by a definite false evaluation:
\begin{equation}
 \operatorname{rej}(C)=\{\nu\in\mathcal N:\exists s\in\nu,\exists c\in C,
   \operatorname{Eval}(c,s)=\mathsf{false}\}.
 \label{eq:main-negative-rejection}
\end{equation}
An unknown evaluation contributes no rejection.  Perturbations form
singleton traces, whereas a cleaned post-exit continuation remains one trace.
The trajectory-level coverage is
\begin{equation}
  \operatorname{cov}(C)=
  \frac{|\operatorname{rej}(C)|}{|\mathcal N|},
  \qquad |\mathcal N|>0.
  \label{eq:method-coverage}
\end{equation}
Each trace contributes one unit regardless of its length.  Instances with no
retained negatives are not used for SFT acceptance; RL uses the binary
all-admitted-clauses-survive fallback defined in
Algorithm~\ref{alg:formal-rl-reward}.

\tightparagraph{Empirical alignment.}
After inductive verification, coverage measures exclusion of sampled
counterfactual behavior.  Candidate negatives are not certified unreachable;
the verifier supplies soundness, while coverage is an empirical strength
signal rather than a guarantee of hidden-target sufficiency.
On 5,711 candidate sets across 691 programs with both positive and negative
final verdicts, coverage achieves a within-program macro AUROC of 0.842
(95\% bootstrap CI: [0.826, 0.858]).  This diagnostic uses archived
target-hidden tool outputs and verified reference sets, supporting an
association with final verification.  Its population and limits are detailed
in Appendix~\ref{app:coverage-predictiveness}.

\subsection{Distilling Multi-Response Search}
\label{sec:method:sft}

SFT aims to distill multi-response compositional search into the policy, so
that a single rollout can reproduce discoveries obtained by combining
multiple responses.  Its target is the composed result, not a selected
standalone response.  It uses Algorithm~\ref{alg:method-pipeline} as a target
synthesizer.  Starting
from an empty accumulated clause pool, each round adds a new group sampled
from the same small open-weight base model that is subsequently fine-tuned.
Executor-derived traces and verifier outcomes jointly determine which
candidate sets become SFT targets.  Each round adds the new base-model clauses
and re-filters the entire accumulated pool:
\begin{equation}
  P_t=\operatorname{Unique}\!\left(P_{t-1}\Vert
    \operatorname{Admit}(y_1^{(t)})\Vert\cdots\Vert
    \operatorname{Admit}(y_{g_{\mathrm{sft}}}^{(t)})\right),
  \qquad
  C_t=\operatorname{Filter}_{\mathcal L}(P_t).
  \label{eq:method-sft}
\end{equation}
Accumulating raw clauses in \(P_t\), rather than only survivors \(C_t\),
allows a removed clause to return when a later response supplies support.
The target can therefore combine discoveries across responses and rounds.

\tightparagraph{Acceptance and serialization.}
Let \(t_{\mathrm{sft}}\) be the synthesis round limit.  For
\(|\mathcal N|>0\), define the first accepted round and its training target by
\begin{align}
 t^*&=\min\{t\in\{1,\ldots,t_{\mathrm{sft}}\}:
             \operatorname{cov}(C_t)\ge\tau\},\nonumber\\
 y^*&=\operatorname{Serialize}(C_{t^*}).
 \label{eq:main-sft-acceptance}
\end{align}
If the set is empty, no training pair is produced; otherwise
\((\mathcal L,y^*)\) enters the SFT corpus without \(\mathcal G\).
The policy is trained with standard autoregressive cross-entropy on the
serialized composition \(y^*\), conditioned on the target-masked loop.
Coverage acceptance does not certify hidden-target sufficiency;
Algorithm~\ref{alg:formal-sft} gives the complete procedure.

\subsection{Full decompose credit}
\label{sec:method:reward}

RL aims to make useful clauses appear earlier in sampling by rewarding the
parts retained after decomposition and joint filtering, even when a complete
response fails.  For a group sampled from the current policy, let
\(A_i=\operatorname{Admit}(y_i)\) as defined above.  We pool and filter once,
then map surviving clauses back to their proposing responses:
\begin{align}
  P^*&=\operatorname{Unique}\!\left(A_1\Vert\cdots\Vert A_{g_{\mathrm{rl}}}\right),
  &C^*&=\operatorname{Filter}_{\mathcal L}(P^*),\\
  V_i&=\{c\in A_i:\operatorname{key}(c)\in\operatorname{keys}(C^*)\},
  &R_i&=\operatorname{rej}(V_i).
  \label{eq:method-credit-pool}
\end{align}
Here \(\operatorname{keys}(C^*)\) is the family of canonical keys of the
pooled survivors.  Thus \(V_i\) assigns each pooled survivor back to every response that proposed
it.  A response retains the coverage contributed by its surviving clauses
even when its other clauses fail, while mutually supporting clauses are
checked in the same pooled context
used at inference.  The cheap reachable-state check is an internal front-end
of this single pooled filtering call; it does not filter responses separately.

The rollout reward is
\begin{equation}
\label{eq:generation-reward}
  r_i=w_c\frac{|R_i|}{|\mathcal N|}.
\end{equation}
Here \(w_c\) is the coverage weight.  The per-response clause cap still
limits admission, but excess clauses incur no reward penalty.

GRPO standardizes rewards within each group and applies each response's
scalar advantage to its tokens.  Thus clauses determine the score, while the
policy update acts on the whole response rather than assigning a separate
policy loss to each clause.  A clause that only supports another clause
receives no direct coverage credit.  The update and empty-negative fallback are detailed
in Appendix~\ref{app:grpo-details} and
Algorithm~\ref{alg:formal-rl-reward}, respectively.

SFT learns to reproduce accepted compositions; RL rewards useful decomposed
parts.  Both use the pooled filtering operation applied at inference, without
access to the hidden target.
